% ****************************************************************************************************
% Chapter 13: Reflections
% ****************************************************************************************************

\chapter{Reflections}
\label{ch:reflections}

This chapter is different in character from the preceding twelve. It contains no new experimental results, no figures, no tables. It is instead a set of reflections on what this project has been as an undertaking, on the methodological choices behind it, on the tools used to carry it out, and on the distance between what a Kaggle competition invites and what a research project becomes when taken seriously. I have written it because I think the methodology of the project is at least as interesting as the findings, and because a report that ends on quantitative results alone would leave unstated some of the most useful lessons I take away from the work.

The chapter is short. It is not a summary of the preceding chapters, those are summarised at the end of each relevant chapter, and the overall picture is consolidated in Chapter~\ref{ch:limitations}'s ``Honest Position.'' What follows here is instead a personal-methodological commentary, organised around four questions: what transferred from my scientific background into this work, what kind of research judgment the project required, how modern language-model assistance shaped the methodology, and what the experience taught me about the difference between competition benchmarking and research-grade machine learning.

\section{Coming to Machine Learning from Another Field}

I trained and worked as an astrophysicist before turning to machine learning. This is my first substantial ML project, and I want to record something about what the transition involved: what transferred, what did not, and what surprised me.

What transferred most readily was the basic discipline of numerical and statistical work: the habit of treating every number as a measurement with an uncertainty, the instinct to visualise data before modelling it, the expectation that noise and signal will coexist in ways that require explicit separation. Astrophysics is a field in which one routinely works with data that are incomplete, censored, and contaminated by systematic effects that nobody has fully characterised. Emergency department data turns out to share many of those properties: missing vital signs at triage, institutional idiosyncrasies in how fields are recorded, patients whose true clinical status can only be inferred from downstream outcomes. The skills for handling that kind of data are recognisably the same.

What transferred less well was the particular form of modelling instinct that physics training cultivates. In astrophysics, one starts with an underlying physical mechanism (gravitational dynamics, radiative transfer, stellar evolution) and derives a model whose parameters have physical meaning. When the fit fails, there is usually a physical hypothesis to test: is the underlying mechanism wrong, is the data calibration off, is there a missing effect? Gradient-boosted tree ensembles do not work like this. There is no underlying mechanism; there is a set of learned decision boundaries whose collective behaviour produces predictions. When the model fails on a subgroup, as mine does on ESI-5, or on young patients with compensated shock, the question is not ``what physics is missing'' but ``what features, training labels, or loss functions would change the boundaries in a useful direction.'' This is a different mode of inquiry, and it took me several months of working with it before the reflex of reaching for a physical mechanism was replaced by the reflex of reaching for a feature-engineering experiment.

What surprised me was how much of the project was about data rather than modelling. The single largest impact on my final performance came from the decision to pivot from the synthetic Kaggle data to MIMIC-IV-ED (Chapter~\ref{ch:discovery}), a data-quality decision, not a modelling decision. The second-largest impact came from feature engineering (Chapter~\ref{ch:features}): the NEWS2 and qSOFA composite scores, the age-stratified abnormal-vitals flags, the augmented features for closing the age disparity. Model selection and hyperparameter tuning, which dominate much of the public discussion of ML, contributed substantially less. This is a lesson the field seems to have learned and forgotten several times: in any real ML problem on tabular data, the data and the features matter much more than the model.

\section{The Decision to Pivot on Real Data}

The most consequential research-judgment decision in the project was the pivot from synthetic Kaggle data to MIMIC-IV-ED. The details of how I discovered the label leakage are in Chapter~\ref{ch:discovery}; what I want to reflect on here is the structure of the decision itself.

The leakage was, in one sense, a gift. It would have been possible to ignore it, to report a text-only F1 of $0.998$ on the Kaggle leaderboard, write up the methodology as if the task had been a real clinical classification problem, and let the synthetic-data artefacts stand as a stand-in for real performance. This would have produced a competitive leaderboard submission. It would also have produced a methodologically dishonest report.

The alternative, which I took, was to treat the discovery of the leakage as grounds for changing what the project was about. Instead of a competition entry that demonstrated high F1 on a compromised benchmark, the project became a demonstration that the benchmark itself was not scientifically usable, and that serious work on the same problem required real clinical data. This was a more expensive path: it required applying for PhysioNet credentialed access, learning a substantially more complex dataset, and building a new pipeline from scratch on real records. It also required accepting that my final F1 on real data ($0.597$) would look much worse than my F1 on the compromised synthetic benchmark ($0.998$), even though the real number represents scientifically meaningful work and the high number represents an artefact.

I think this kind of decision is the most important thing a research-oriented practitioner can learn to make. Competitions reward specific performance metrics on specific benchmarks. Research rewards methodological honesty, generalisability, and clinical relevance. When these diverge, and they often do, choosing research over competition is not a moral stance; it is simply the only way to produce work that anyone should act on. The compromised leaderboard submission would have been a win on the Kaggle metric and a null result for the underlying problem. The pivot produced work that actually addresses the underlying problem, at the cost of looking less dominant on the immediate competition.

A broader observation follows from this. I am not the first team to find label leakage in a synthetic ML benchmark. I will not be the last. As synthetic data generation becomes more common, the skill of identifying when a benchmark is not measuring what it claims to measure will become more important, not less. Some of the most valuable contributions a researcher can make are negative: this benchmark is not scientifically usable, this result does not generalise, this improvement is an artefact. I wish the reward structures of the field supported this kind of work more consistently than they do.

\section{LLM-Assisted Research as a Methodology}

This project was carried out with substantial assistance from large language models, primarily Claude (Anthropic). I want to describe how that assistance was used, because the methodology is still novel enough that practitioners will benefit from concrete examples of what works and what does not.

The assistance operated at three distinct levels. At the lowest level, language models served as an enhanced form of pair programming: I described what I wanted to compute, the LLM produced a first draft of code, I ran it, debugged errors together, and iterated until the code produced the intended output. This use is continuous with decades of autocomplete and code-generation tools, though the contextual reasoning of modern models makes the loop substantially tighter.

At the middle level, language models served as a research assistant for literature review and methodological guidance. I could describe my task (five-class ordinal classification with severe class imbalance and asymmetric clinical costs) and ask for summaries of the relevant literature (Mirhaghi's meta-analysis, the RETTS validation work, the various Kaggle ED-triage papers) with references. The summaries were not always accurate on specifics, and every citation was independently verified before inclusion in this report. But the models were quite good at sketching the rough topology of a field I was new to, suggesting methodological angles (focal loss, asymmetric weights, ordinal penalties) that I might not have found on my own.

At the highest level, and the one that matters most, language models served as a collaborator in scientific writing. Many of the paragraphs in this report were first drafted by me, then edited by a language model for clarity or structure, then revised again by me. Some paragraphs were first drafted by a language model given a detailed outline and source material, then heavily edited by me. The final text is thoroughly my own in the sense that every sentence has been read, considered, and accepted by me, but it would be dishonest to claim that the prose is purely human in origin.

The methodological discipline that makes this work is \emph{verification}. Language models are fluent, plausible, and often wrong. They produce confident-sounding prose about numerical results that they have not actually computed, citation suggestions for papers that do not exist, and explanations of mechanisms that sound right but fail when checked against the data. The only way to use them safely in scientific work is to verify everything: every number must trace to a saved computational output, every citation must be retrieved and read, every claim about what the model does must be checked against the actual model.

I developed a specific workflow for this. Computational results were computed once by scripts, saved to disk as structured files (CSVs, JSON summaries, NumPy arrays), and then separately verified by a second process that read the saved files and produced a verification report. The verification reports, eight of them for this project, one for each major experimental result, were used as the authoritative source for every number in this report. The language model assisting with writing was given the verification reports, not the raw computational logs. This layer of separation prevented several potential errors where the model might have confabulated numbers that were not actually present in the outputs. In at least one case (the ``calibration improves the operating curve'' claim corrected in Chapter~\ref{ch:calibration}), the verification process caught an incorrect statement in an earlier summary and forced me to reframe the finding honestly.

I believe this is the right methodology for LLM-assisted research: heavy use of the tools for drafting, literature work, and code generation; disciplined separation between computational outputs and writing support; verification as a mandatory layer between raw results and final claims. Without verification, LLM-assisted research is not research at all, it is plausibly-worded fiction. With verification, the productivity gains are substantial and the outputs can be held to ordinary scientific standards.

A final note on acknowledgment. I describe this collaboration openly because the alternative, pretending the prose was written without assistance, would be misleading about the methodology of modern research. A reader evaluating the report is entitled to know that the text was produced by a specific kind of collaboration, and that the reliability of its claims rests on the verification discipline described above. I suspect this kind of openness will become standard in research practice over the coming years. For now, it is still unusual enough to be worth explaining.

\section{Competition and Research}

The project was formally a Kaggle competition entry. It became, in practice, a small research project with a Kaggle writeup attached. The gap between those two framings deserves a brief comment.

Competitions reward specific performance on specific benchmarks. The rubric for this competition allocates points for data exploration, modelling, bias and fairness, explainability, calibration, and documentation; teams are rewarded for hitting each category, and the writeup is graded on evidence that each was addressed. This is a reasonable structure for a competition; it is not a reasonable structure for research.

A research project on the same topic would allocate effort differently. It would spend much more time on external validation (which this competition does not require), on prospective outcome linkage (which neither competition nor my time budget allowed), on clinician-interaction studies (which require deployment infrastructure I do not have), and on multi-institution generalisation (which requires additional credentialed data access). It would spend less time on, say, documentation polish or on making the writeup legible to a non-technical reviewer.

I have tried to take the best of both framings. From the competition side, I took seriously the rubric's insistence on fairness auditing, explainability, and calibration, these are categories that real research projects sometimes neglect and that the rubric pushed me to address. From the research side, I took seriously the obligation to report null results (calibration does not improve the operating curve at matched multipliers; asymmetric weights and per-class thresholds are substitutes, not complements; the SHAP-derived explainability has an important interpretability gap); to disclose limitations prominently rather than bury them; and to position the work accurately relative to the underlying clinical problem. The result is a document that is more than a leaderboard submission but less than a peer-reviewed paper, a research demonstration with the structure of a competition writeup.

Whether this hybrid form is useful to the field is not for me to judge. I have written it in the hope that it is.

\section{What the Project Has Taught Me}

A few lessons seem clear enough to state directly.

First, the most important findings of this project were not the positive ones, not the F1 numbers, not the Cohen's kappa, not the feature-importance rankings. The most important findings were a series of disappointments that turned out to be informative. The synthetic benchmark did not measure what it claimed to measure. Calibration did not improve the operating curve the way I initially claimed. Asymmetric training weights did not complement threshold optimisation, they duplicated it. Race was not the primary axis of my model's bias, age was. Each of these findings would have been a negative result by the standards of a competition leaderboard. Each of them is scientifically more useful than a positive result would have been.

Second, the hardest part of the project was not the modelling. The hardest parts were the data-quality investigation that led to the pivot, the careful selection of features that would encode what a clinician actually looks at, the audit that turned up the age disparity, and the verification discipline that prevented me from publishing wrong numbers. Modelling, in the sense of ``choose a model architecture and tune hyperparameters,'' was close to routine.

Third, domain knowledge matters more than I would have guessed before starting. The clinical literature on compensated shock (Chapter~\ref{ch:fairness}), on the ESI scoring system's known limitations (Chapter~\ref{ch:triage-problem}), on the Mirhaghi meta-analysis of triage reliability (Chapter~\ref{ch:results}), and on the EU AI Act's specific compliance provisions (Chapter~\ref{ch:explainability}) all shaped my methodology in ways that a purely technical ML team would have missed. This is not a surprise to anyone who has worked at the intersection of ML and a specific application domain, but it is perhaps a surprise to a practitioner coming from a more quantitative field where the domain and the method are often the same thing.

Finally, I close with an observation about the place of machine learning in clinical work. My model matches human triage reliability; it does not exceed it. At most, what it offers is a second opinion: a consistent, always-available comparison that can catch cases where the human decision drifts. This is not a modest contribution, Mirhaghi's meta-analysis suggests that human triage is unreliable enough that a consistent second opinion might produce real benefits, but it is a very different contribution from the ``ML replaces clinician'' framing that has sometimes been sold. The appropriate deployment of a system like mine is as a decision-support tool within a workflow that retains full human authority. Everything in the preceding twelve chapters has been designed to support that framing rather than to contradict it.

The references follow. I am grateful to the Laitinen--Fredriksson Foundation for convening this competition, to PhysioNet and the MIMIC-IV-ED maintainers for the data access that made the real work possible, and to the authors of the many open-source tools, LightGBM, XGBoost, CatBoost, scikit-learn, SHAP, Hugging Face Transformers, and the rest, that my pipeline rests on.
